\section{Background and related work}\label{sec:background}

\subsection{Selecting the number of clusters}\label{subsec:bg-k}

Choosing the number of clusters $K$ is a classical problem in unsupervised
analysis, and a large family of \emph{internal} criteria address it by scoring
a partition against the geometry of the data alone. The gap statistic
\citep{Tibshirani2001} compares the within-cluster dispersion to its
expectation under a reference distribution with no clustering, typically
uniform over the data's bounding box; the average silhouette
\citep{Rousseeuw1987}, the Calinski--Harabasz ratio
\citep{CalinskiHarabasz1974} and the Davies--Bouldin index
\citep{DaviesBouldin1979} balance within- against between-cluster scatter.
\revtext{Comparative studies cover the many further variants
\citep{Arbelaitz2013}. A second family selects $K$ through a probabilistic
model: model-based clustering fits a finite mixture and chooses the number of
components by an information criterion such as the BIC
\citep{GormleyMurphyRaftery2023}. None of these criteria was proposed with a
distribution theory that requires independent observations. What they share is
that the dependence between observations plays no role in their construction:
the reference distribution of the gap statistic, the implicit baselines of the
ratio indices, and the mixture likelihood treat the sample as exchangeable
draws, and so ignore spatial dependence rather than model it. When
observations are nearly independent the criteria are well calibrated; under
strong dependence they are not.} Spatial data are the archetypal failure,
because nearby
locations are correlated: ``everything is related to everything else, but
near things are more related than distant things'' \citep{Tobler1970}.

\subsection{Accounting for spatial dependence}\label{subsec:bg-spatial}

That an independence reference over-detects structure in dependent data is
known across several fields. In spatial epidemiology, \citet{Goovaerts2004}
show that a complete-spatial-randomness null over-identifies disease clusters
and propose, in its place, a geostatistical \emph{spatial neutral model}
(realisations reproducing the histogram and variogram of the data by sequential
Gaussian simulation), against which far fewer locations remain significant.
In neuroimaging, \citet{Eklund2016} trace inflated false-positive rates in
cluster-extent inference to a mis-specified parametric model of the spatial
autocorrelation. The common remedy is a null that \emph{preserves} the
spatial covariance structure while removing the effect being tested, realised
either by geostatistical simulation \citep{Goovaerts2004} or by spatially
constrained randomisation such as Moran spectral randomization
\citep{WagnerDray2015}.

The same logic has been brought to bear on $K$ itself.
\citet{HennigLin2015} calibrate a validity index (the gap, silhouette or
prediction strength) against its distribution under a problem-specific null
that encodes non-clustering structure, including spatial autocorrelation, and
use the calibrated index both to test for clustering and to estimate $K$; on a
biogeographic example a plain Gaussian null favours eight clusters while the
autocorrelation-aware null reduces this to two and renders the apparent
clustering non-significant. Their framework is the conceptual basis of the
present work. \revtext{Its guiding principle, that all structure not
attributable to clustering belongs in the null model, is exactly the principle
the present method instantiates for spatially autocorrelated continuous
fields; the two contributions differ in how the calibration is obtained and in
the data setting they treat. First, \citet{HennigLin2015} calibrate by
simulation, so every decision requires an ensemble of null datasets, whereas
the calibration developed here is analytic, through the effective sample size,
which is what makes the near-linear full-scene cost of
Section~\ref{subsec:complexity} possible. Second, their construction} is
formulated for discrete presence--absence data
with an algorithmic disjunction-parameter null, and addresses neither the
\emph{continuous-field} setting of gridded imagery nor the confounding of a
smooth large-scale \emph{trend} with discrete classes. The present
contribution supplies exactly these: a Gaussian-random-field null on the
lattice, a within-class range obtained by local range estimation, an adaptive
trend--residual separation, and an effective-sample-size calibration of a
modality statistic. \revtext{Throughout the comparisons of
Sections~\ref{sec:simulation} and~\ref{sec:application}, the label
Hennig--Lin accordingly denotes their published null-model construction
applied to our setting, not this general framework, of which the proposed
method is itself an instance.}

The quantitative device that makes this possible is the \emph{effective
sample size}: under spatial autocorrelation the variance of a sample statistic
behaves as if computed from $\neff<n$ independent observations, a reduction
governed by the integral of the autocorrelation \citep{Cressie1993}\revtext{,
a quantity formalised as the \emph{integral range} in the theory of ergodic
random functions \citep{Lantuejoul1991,Lantuejoul2002}}.
\revtext{The concept recurs across statistics whenever dependent data carry
less information than their count suggests: \citet{Griffith2005} develops
effective geographic sample sizes for spatial sampling design,
\citet{Berger2014} formalise the notion for general parametric inference,
\citet{VallejosOsorio2014} define and study it for spatial process models, and
\citet{Acosta2021} extend its computation to large spatial datasets through
block likelihoods.}
\citet{Dutilleul1993} operationalised this idea to correct the degrees of
freedom of a correlation test between two spatial processes; we use it to
recalibrate, rather than to correct after the fact, the reference distribution
of a clustering criterion. The autocorrelation scale it requires is the
variogram range \citep{Matheron1963,Cressie1993}, and the safeguard against
letting large homogeneous regions dominate its estimate is \emph{declustering},
a standard geostatistical correction for spatially clustered support
\citep{IsaaksSrivastava1989,Deutsch1998}.

\subsection{Modality as a clustering signal}\label{subsec:bg-modality}

Because $k$-means dispersion is itself a variance-based quantity, a null
matched to the data's covariance structure can explain away the very clusters one
seeks; a criterion that survives such a null must use information beyond the
covariance. Multimodality is the natural choice. The dip test of
\citet{HartiganHartigan1985} measures the departure of a univariate sample from
unimodality, and dip-means \citep{Kalogeratos2012} turns it into a
divisive procedure that splits a cluster only when a projection of it is
significantly multimodal, thereby estimating $K$. These methods assume
independent observations; our contribution is to make the dip test valid under
spatial autocorrelation by calibrating it at the effective sample size, and to
pair it with the trend removal that prevents a smooth gradient from being read
as two modes.

\subsection{Distinction from spatially constrained clustering}\label{subsec:bg-constrained}

A separate line of work \emph{imposes} spatial contiguity on the partition:
regionalization methods such as SKATER \citep{Assuncao2006}, ClustGeo
\citep{Chavent2018} and the max-$p$-regions problem \citep{Duque2012} constrain
clusters to be spatially connected. These address a different question, how
to partition space into contiguous regions, and still require the number of
regions as an input or a tuning parameter. We do not constrain the partition;
we correct the criterion that selects \emph{how many} groups are present, and
the two are complementary: the effective-sample-size calibration developed here
could in principle inform the stopping rule of a regionalization just as it
informs the recursion of Section~\ref{sec:method}.
